\begin{frame}{왜 정규분포인가 --- 세 가지 이유가 겹침}
  \begin{enumerate}
    \item \kb{지지가 전체} $\mathbb{R}^d$ --- 어떤 연속 행동도
      이론적으로 가능한 확률을 가짐.
    \item \kb{미분이 깨끗하다} --- 평균에 대해선
      $\nabla_\mu \log\pi = \Sigma^{-1}(a-\mu)$ 로 \textbf{선형},
      자유도 $d(d+1)/2$ 개.
    \item \kb{최대 엔트로피 원리} --- 이산의 softmax가 합 1 제약하의
      ``최대 엔트로피 분포''라면, 연속에서는 같은 원리가
      \textbf{정규분포}를 만듦. (아래 FAQ에서 증명)
  \end{enumerate}
  \vskip0.5em
  \begin{block}{한 줄}
    세 이유가 겹치므로 ``연속 정책의 표준 파라미터화 = 조건부 정규분포''.
  \end{block}
\end{frame}
